Whether the field moves a dislocation was checked directly. An aluminium slab
contains the same held inclusion and a straight edge dislocation on the slip plane
12~\AA{} above it (Supplementary S6). With the inclusion unstrained, the
dislocation lies, after energy minimisation, «DISL_D»~\AA{} from the inclusion centre along $x$,
where the line-averaged shear stress of the field is about «DISL_TAU»~MPa. When the inclusion is given the strain~(\ref{eq:strain}), the
dislocation moves «DISL_SHIFT»~\AA{} away from it, about half the distance
$b/2=1.43$~\AA{} between neighbouring valleys of the lattice relief (the
positions where the energy of the dislocation is lowest), and does not
pass into the next valley. When the strain is removed, the dislocation
returns to within «DISL_REV»~\AA{} of where it was, so in this calculation at
0~K the response is elastic and reversible. The line-averaged field on the
dislocation is about half the lattice resistance of the potential, 9~MPa
(Supplementary S6). Where the resistance is lower, for example at
room temperature, the same field could move the dislocation further; in the
alloy, Mg and Si atoms and precipitates also hold dislocations, and the
resistance there is higher (13--20~MPa).
